\begingroup
\renewcommand{\arraystretch}{1.45}
\fancyhead[L]{\small BESTIMMUNG DER ANZIEHUNG}
\Needspace{30\baselineskip}
\AttractionArticle{12}

Nachdem wir gezeigt haben, wie sich die zwölf Größen $G$, $G'$, $G''$, $\alpha$, $\alpha'$, $\alpha''$, $\beta$. $\beta'$, $\beta''$, $\gamma$, $\gamma'$, $\gamma''$ vollständig bestimmen lassen, gehen wir zur Entwicklung des Differentials $\mathrm dE$ über. Setzen wir
\[
\begin{aligned}
&t=\gamma+\gamma'\cos T+\gamma''\sin T\quad
\mathinner{\ldotp\;\ldotp\;\ldotp\;\ldotp\;\ldotp\;\ldotp\;\ldotp\;\ldotp}\quad[20]
\end{aligned}
\]
sodass gilt
\[
\begin{aligned}
&t\cos E=\alpha+\alpha'\cos T+\alpha''\sin T\quad
\mathinner{\ldotp\;\ldotp\;\ldotp\;\ldotp\;\ldotp\;\ldotp\;\ldotp}\quad[21]
\end{aligned}
\]
\[
\begin{aligned}
&t\sin E=\beta+\beta'\cos T+\beta''\sin T\quad
\mathinner{\ldotp\;\ldotp\;\ldotp\;\ldotp\;\ldotp\;\ldotp\;\ldotp}\quad[22]
\end{aligned}
\]
\Needspace{27\baselineskip}
Hieraus leiten wir ab
\[
\begin{aligned}
&t\,\mathrm dE=\cos E\,\mathrm d.\,t\sin E-\sin E\,\mathrm d.\,t\cos E\\[8pt]
&=\cos E(\beta''\cos T-\beta'\sin T)\mathrm dT
-\sin E(\alpha''\cos T-\alpha'\sin T)\mathrm dT
\end{aligned}
\]
und daher
\[
\begin{aligned}
&tt\,\mathrm dE=(\alpha\beta''-\alpha''\beta)\cos T\,\mathrm dT
+(\alpha'\beta-\beta'\alpha)\sin T\,\mathrm dT
\\[6pt]
&\qquad{}+(\alpha'\beta''-\beta'\alpha'')\mathrm dT\\[8pt]
&=\varepsilon\gamma'\cos T\,\mathrm dT
+\varepsilon\gamma''\sin T\,\mathrm dT
+\varepsilon\gamma\,\mathrm dT=\varepsilon t\,\mathrm dT
\end{aligned}
\]
oder
\[
\begin{aligned}
&t\,\mathrm dE{}=\varepsilon\,\mathrm dT\quad
\mathinner{\ldotp\;\ldotp\;\ldotp\;\ldotp\;\ldotp\;\ldotp\;\ldotp\;\ldotp\;\ldotp\;\ldotp\;\ldotp\;\ldotp\;\ldotp\;\ldotp\;\ldotp\;\ldotp\;\ldotp\;\ldotp\;\ldotp\;\ldotp\;\ldotp}\quad[23]
\end{aligned}
\]
\Needspace{18\baselineskip}
\par Es ist zu beachten, dass die Größe $t$ ihrer Natur nach stets positiv ist, wenn der Koeffizient $\gamma$ positiv ist, oder stets negativ, wenn $\gamma$ negativ ist. Da nämlich
$(\gamma'\cos T+\gamma''\sin T)^2+(\gamma''\cos T-\gamma'\sin T)^2=\gamma'\gamma'+\gamma''\gamma''=\gamma\gamma-1$ gilt, wird stets
\noindent\SourceLeaf{358}
$\gamma'\cos T+\gamma''\sin T$, ohne Rücksicht auf das Vorzeichen, kleiner als $\gamma$ sein. Hieraus schließen wir, dass, sooft $\varepsilon\gamma$ eine positive Größe ist, die Variablen $E$ und $T$ stets zusammen zunehmen; ist aber $\varepsilon\gamma$ eine negative Größe, so muss die eine Variable stets abnehmen, während die andere zunimmt.

\Needspace{28\baselineskip}
\AttractionArticle{13}

Der Zusammenhang zwischen den Variablen $E$ und $T$ wird durch die folgenden Überlegungen noch deutlicher. Setzt man $\surd{(\gamma\gamma-1)}=\delta$, sodass $\delta\delta=\alpha\alpha+\beta\beta=\gamma'\gamma'+\gamma''\gamma''$ gilt, so leiten wir aus den Gleichungen 20, 21, 22 ab
\[
\begin{aligned}
&t(\delta+\alpha\cos E+\beta\sin E)\\[8pt]
&\quad=\gamma\delta+\alpha\alpha+\beta\beta
+(\gamma'\delta+\alpha\alpha'+\beta\beta')\cos T
+(\gamma''\delta+\alpha\alpha''+\beta\beta'')\sin T\\[8pt]
&\quad=(\gamma+\delta)(\delta+\gamma'\cos T+\gamma''\sin T)
\end{aligned}
\]
\Needspace{30\baselineskip}
\par Ebenso folgt aus den Gleichungen 21, 22
\[
\begin{aligned}
&t(\alpha\sin E-\beta\cos E)=\varepsilon(\gamma'\sin T-\gamma''\cos T)
\end{aligned}
\]
Setzt man
\[
\begin{aligned}
&\frac{\alpha}{\delta}=\cos L,\quad
\frac{\beta}{\delta}=\sin L,\quad
\frac{\gamma'}{\delta}=\cos M,\quad
\frac{\gamma''}{\delta}=\sin M
\end{aligned}
\]
so erhalten diese Gleichungen die folgende Form:
\[
\begin{aligned}
&t(1+\cos(E-L))=(\gamma+\delta)(1+\cos(T-M))
\end{aligned}
\]
\[
\begin{aligned}
&t\sin(E-L)=\varepsilon\sin(T-M)
\end{aligned}
\]
woraus sich durch Division wegen $(\gamma+\delta)(\gamma-\delta)=1$ ergibt
\[
\begin{aligned}
&\mathrm{tang}\,\frac12(E-L)=\varepsilon(\gamma-\delta)\mathrm{tang}\,\frac12(T-M)
\end{aligned}
\]
\[
\begin{aligned}
&\mathrm{tang}\,\frac12(T-M)=\varepsilon(\gamma+\delta)\mathrm{tang}\,\frac12(E-L)
\end{aligned}
\]
\Needspace{18\baselineskip}
\par Hieraus ergibt sich nicht nur dieselbe Folgerung, zu der wir am Ende des vorangehenden Artikels gelangt sind; außerdem ist auch ersichtlich, dass, wenn der Wert von $E$ um 360 Grad zunimmt, der Wert von $T$ um ebenso viel entweder zu- oder abnimmt, je nachdem, ob $\varepsilon\gamma$ eine positive oder eine negative Größe ist. Setzt man ferner $\delta=\mathrm{tang}\,N$, $\gamma=\sec N$, so wird offenbar gelten
\[
\begin{aligned}
&\gamma-\delta=\mathrm{tang}(45^0-\frac12N),\quad
\gamma+\delta=\mathrm{tang}(45^0+\frac12N)
\end{aligned}
\]
\SourceLeaf{359}
\clearpage
\endgroup
